\documentclass[12pt]{article}
\usepackage[utf8]{inputenc}
\usepackage[margin=0.75in]{geometry}
\usepackage{amsmath}
\usepackage{amssymb}
\usepackage{amsthm}
\usepackage{graphicx} % Required for inserting images
\usepackage{array}

\title{Norm Emergence Report: Social Reward Economies}
\author{Aaron Avram}
\date{September 24th 2026}

\begin{document}
\theoremstyle{plain}
\newtheorem{theorem}{Theorem}
\newtheorem{proposition}{Proposition}
\newtheorem{lemma}[theorem]{Lemma}
\newtheorem{corollary}[theorem]{Corollary}
\newtheorem{claim}[theorem]{Claim}

\theoremstyle{definition}
\newtheorem{definition}{Definition}
\newtheorem{example}{Example}
\newtheorem{exercise}[theorem]{Exercise}

\theoremstyle{remark}
\newtheorem{remark}[theorem]{Remark}
\newtheorem{note}[theorem]{Note}

\maketitle

\section{Introduction}
In this report, I wish to develop a model for a multi-agent system with two social reward economies
within the frameworked proposed in the \textit{Social Reward Economies} paper. We use multiple real-life
examples as guides to the right formal definitions, with the \#MeToo movement as a primary case study.

A social reward economy is defined formally in Section 3, but intuitively
it is a network of socially responsive agents that want to maximize their perceived social reward.
They do this by choosing schema and behaviours that produce positive group feedback while while avoiding
negative feedback.

This report only considers the case where a group of agents $G$ in a social reward economy contains a subgroup $G_r$
which has its on social reward economy. The norm of the former economy can be thought of as the "dominant norm"
with the latter being the "resistant norm". These two groups and their social reward economies will be defined
precisely in Sections 4 and 5 respectively. A major point of interest will be how these two groups
and their differing norms interact, which will be explored in detail in Section 6.

\section{Modelling the Real World}
First we explore some real-world examples to better understand
how to motivate the formal definitions.

A first example is the formation of the \#MeToo Movement. In this situation, society
could be taken as the dominant group, and young and middle aged women as the resistant subgroup.
These women are unhappy with the dominant norm, which rewards keeping quiet about
sexual harassment in the workplace. As a subgroup, they instead reward speaking
up about sexual harassment and shaming the individuals who do these things.
Hence, the reputation of abusers in the resistant group will be highly negative
but in the dominant group their reputation would be average or even potentially very high
if they occupy a position of power (e.g. a high ranking corporate executive). Conversely,
the reputation of a member of the \#MeToo movement in the movement itself will be quite high,
but in the dominant group it would be lower or potentially negative, depending on how strongly
the movement has caught on within the larger group.

Here it is important to actually distinguish between three types of people:
those in charge in the dominant group (such as high ranking corporate officials)
the members of the \#MeToo movement and finally the general public (which is just everyone else in society). In this setup
it is clear that those in charge benefit at the expense of many women in the workplace
and by extension many members of the \#MeToo movement. Hence, the group in charge
will have an incentive to preserve the dominant norm, which advocates for not calling out sexual harassment,
as this directly benefits them. Conversely, women and the \#MeToo movement wish to overturn the dominant
norm in favour of their resistant one, that sexual harassment should be made public and the perpetrators
should be shamed, as this benefits them. However, the general public is not strongly affected by the behaviours
of either of the above groups, yet they make up a majority of the population. Hence
both of the previous groups aim to enforce their respective norm by influencing the general public
to believe in their viewpoint. Members of the \#MeToo movement attend rallies, post on social media
and do other things to try and spread awareness to the general public of their unjust treatment. Meanwhile,
those in charge denounce the accusations of the \#MeToo movement and try to convince
the general public the people involved in this movement are not credible and should not be taken seriously.
The norm that succeeds from this will be the one which the general public becomes most willing to accept.

A second example is a situation where the dominant group does not necessarily
benefit at the expense of the subgroup. An example of this could be the
movement advocating for improved accessibility on public transit. The general public and those in charge
do not gain anything from the lack of accessibility, but those with physical impairments
do suffer disproportionately. Hence, this movement will have a high chance of convincing the rest
of the individuals to update their norm to support accessibility on public transit.

A third example is sports fans. Consider the city of Liverpool as the dominant group,
then there is a subgroup consisting of Liverpool football club fans. This subgroup
has its norm, with accepted behaviours specific to their group (such as going to watch every Sunday match).
Additionally, this subgroup has its own network of communication that is separate from that of the general public.
For instance, many fans may discuss the current state of the club on team-specific online forums that
are not used by the general public. It is important to note that the general public and the fans
are practically unaffected by each other. Hence, these two groups can coexist together with
the fans as a subculture.

With these motivating examples, the objects of study can be defined precisely.

\section{Social Reward Economies}
The definition of an individual social reward economy borrows largely from the revised report from last year.
Given an ambient set of agents $G = \{ 1, \ldots, n \}$, forming a multi-agent system,
a social reward economy with respect to $G$ consists of:
\begin{itemize}
    \item \textbf{The Agents:} $\mathcal A \subseteq G$ this is the set of agents who actively participate in the social reward economy.
    \item \textbf{Social situations:} there is a fixed set of states $S$ and at each timestep the system is in a given state $s \in S$,
    where $p(s)$ is the probability of the system being in state $s$ (this probability is independent of time).
    \item \textbf{Actions:} for each state $s$ there is a finite set of actions $A(s)$ from which
    each agent $i \in \mathcal A$ must choose in state $s$.
    \item \textbf{Behaviour:} the behaviour of an agent $i$ is given by a mapping
    \[ \pi_i: S \to \bigcup_{s \in S}A(s) \]
    such that $\pi_i(s) \in A(s)$.
    \item \textbf{Schemas:} each agent interprets a given state $s$ through an internal schema.
    Let $\Sigma(s)$ denote the set of possible schemas applicable to situation $s$. Schemas
    are not observable, but the behaviour of agent serves as a proxy for their schema. In other words
    when agent $i$ selects a schema $\sigma \in \Sigma(s)$ their corresponding action $\pi_i(s) \approx \sigma$.
    \item \textbf{Utilities:} For each state $s$ and action $a \in A(s)$ the agent $i$ recieves a benefit $u_i(s, a)$.
    \item \textbf{Reward Structure:} there are two core reward systems that serve as the incentive mechansisms for agents:
    \begin{enumerate}
        \item[(1)] Personal utility: this is simply the personal benefit described in the bullet point above.
        \item[(2)] Reputation: the reputation of agent $i$ is defined as the aggregate benefit it provides to other agents.
        If the system is in state $s$ then
        \[ R_i(s) = \gamma \sum_{j \in \mathcal A \setminus \{ 0 \}} \sum_{s \in S}u_j(s, \pi_i(s)) \]
        where $\gamma > 0$ is a scaling parameter.
    \end{enumerate}

    \item \textbf{Reputation Learning through Gossip:} Each agent $i$ maintains two estimates about another agent $j$.
    \begin{enumerate}
        \item[(1)] $P^{(t)}_{i, j}$ denotes the $i$'s personal estimate of $j$'s utility for $i$.
        \[ P^{(t)}_{i, j} = \begin{cases}
                u^{(0)}_{j}(s, \pi_i(s)) & t = 0 \\
                (1 - \alpha)P_{i, j}^{(t-1)} + \alpha u_j^{(t)}(s, \pi_i(s)) & t > 0
        \end{cases} \]
        \item[(2)] $R^{(t)}_{i, j}$ the reputation estimate of $i$ for $j$
        \[ R^{(t)}_{i, j} = \begin{cases}
            0 & t = 0 \\
            R_{i, j}^{t - 1} + (P_{i, j}^{(t)} - P_{i, j}^{(t-1)}) & t > 0
        \end{cases} \]
        finally, after these updates the set of agents $\mathcal A_{a}(t) \subseteq \mathcal A$
        that are active at time $t$ exchange their reputation estimates
        \[ R_{i, j}^{(t)} = \frac{1}{|\mathcal A_{a}(t)}\sum_{k \in \mathcal A_a(t)}R^{(t)}_{k, j} \]
        and this is for all $i \in \mathcal A_{a}(t)$ and $j \in \mathcal A$.
    \end{enumerate}

    \item \textbf{Optimization:} Opinion leaders maximize reputation through a loop
    of behaviour and feedback, other agents imitate the agent who they estimate to have the highest
    reputation.
\end{itemize}

\section{The Dominant Group}
Now we define the dominant group in isolation with the social reward economy framework.

\begin{itemize}
    \item \textbf{The Agents:} $\mathcal A = G$, in the above examples the dominant group consisted
    of the entire society, so it makes sense for its economy to involve all agents.
    \item \textbf{Social situations:} we fix a set of states $S$ which the members of $G$ act in,
    this should be thought of as a large set of states (the different parts of functioning in a society).
    However there is a subset $S_r \subseteq S$ which is the set of states where those in power benefit at
    the expense of members of the \#MeToo movement and who they stand for. 
    \item \textbf{Actions:} we fix the set of actions $A(s)$ for each state $s$.
    \item \textbf{Behaviour:} This is defined as above.
    \item \textbf{Schemas:} In our \#MeToo examples there were roughly 3 schemas. The dominant schema $\sigma_d$
    enforced by those in power, the reformed schema $\sigma_f$ that the \#MeToo movement wants to establish as the norm,
    and finally the resistant schema $\sigma_r$ which is the schema employed by the \#MeToo movement to overturn the dominant one.
    Note that the last two can be folded into a single subgroup schema, which is a better framework to describe
    the accessibility movement or sports fans in relation to society.
    \item \textbf{Utilities:} clearly the distribution of the utilities dictates how the two groups interact.
    We only consider the \#MeToo movement example in detail, but hopefully it is clear how the below utility distribution
    would inform the other examples. For the \#MeToo movement example, we can partition $G$
    as $G_d \sqcup G_n \sqcup G_r$ where $G_d$ are those in power, $G_n$ is the general public
    and $G_r$ are members of \#MeToo. And if we assume all actions fall under one of the 3
    schemas $\sigma_d, \sigma_r, \sigma_f$ we can map out the reward distribution.
    First for $s \in S \setminus S_r$ the rewards are fairly uniform, there is no unfair
    advantage for those in power. However for $s \in S_r$ we have the following distributions where $\gg$ refers to large reward
    and $>$ refers to a normal level of reward:
    If $i \in G_d$
    \begin{align*}
        u_i(s, a) = \begin{cases}
            \gg 0 & a \approx \sigma_d \\
            < 0 & a \approx \sigma_r \\
            > 0 & a \approx \sigma_f
        \end{cases}
    \end{align*}
    the intuition being that those in power benefit greatly from their schema,
    but during a successful resistance they would suffer due to public shaming and loss
    of credibility. However, once the resistance has overturned the dominant schema,
    those in power can reintegrate into society but with moderate rewards and lower reputation
    due to the negative image they have accumulated.
    If $i \in G_n$
    \begin{align*}
        u_i(s, a) = \begin{cases}
            > 0 & a \approx \sigma_d \\
            > 0 & a \approx \sigma_r \\
            > 0 & a \approx \sigma_f
        \end{cases}
    \end{align*}
    here the idea is that the general public is mostly unaffected by the conflict
    between those in power and the \#MeToo movement.
    Finally, if $i \in G_r$
    \begin{align*}
        u_i(s, a) = \begin{cases}
            \ll 0 & a \approx \sigma_d \\
            \gg 0 & a \approx \sigma_r \\
            > 0 & a \approx \sigma_f
        \end{cases}
    \end{align*}
    her the intuition is that, in the dominant norm the individuals the \#MeToo movement
    stand for are negatively affected. Thus, during a successful resistance they would
    see large benefits as their suffering is acknowledged and those who caused harm to them are punished.
    Finally, once the resistant norm becomes the dominant one, those who once suffered will
    now function as normal members of society with moderate rewards in these social situations.
    \item \textbf{Reward Structure:} the two quantities here are defined in the exact same way as above.

    \item \textbf{Reputation Learning through Gossip:} these estimates are defined in the same way as above,
    although it is worth commenting briefly on the gossip network. Since the dominant group will have higher rewards
    on average than both the general public and those in the \#MeToo movement, they will act more often in the system
    and hence their reputation estimates will disproportionately influence the reputation
    estimates of other agents in the dominant norm. This is how those in power aim to influence the general
    public to discredit the \#MeToo movement and continue to follow the dominant norm.

    \item \textbf{Optimization:} the optimization process is identical to before.
\end{itemize}

\section{The Resistant Group}
Now we define the resistant group in isolation with the social reward economies framework.
\begin{itemize}
    \item \textbf{The Agents:} $\mathcal A = G_r$, clearly the agents involved in the resistant economy
    will only be those in the resistant subgroup.
    \item \textbf{Social Situations:} we consider $S_r \subseteq S$ to be the set of states within this social reward economy.
    We only take a subset of the total states in the dominant economy as this provides better flexibility in modeling real-life
    resistance. For example, the \#MeToo movement is only concerned with a small subset of the total states in society,
    namely those where sexual harassment takes place. Thus it would make sense for the social interactions of the group
    to revolve primarily around these specific situations.
    \item \textbf{Actions:} we fix the same set of actions $A(s)$ used in the dominant economy for each state $s \in S_r$.
    \item \textbf{Behaviour:} this is also defined as above, except now in the resistant economy,
    actors are only interested in the actions taken in the subset of states $S_r$.
    \item \textbf{Schemas:} the same schemas as in the previous economy are used here.
    \item \textbf{Utilities:} the utilities are inherited from the dominant economy.
    \item \textbf{Reward Structure:} the two quantities are defined as above, except
    now reputation is taken over $S_r$.
    \item \textbf{Reputation Learning through Gossip:} these estimates are defined in the same way as above,
    but now the gossip network functions slightly differently. Since the only states in the economy
    are those in $S_r$ and only agents in $G_r$ participate in the economy, the emergent opinion leader
    would be the individual who best represents the overall interests of the group as the utility distribution
    of all agents in the group are similar. In the example, this would be the leader of the \#MeToo movement.
    This individual is the candidate reformed opinion leader, at least temporarily.
    In other words, to overturn the incumbent norm the resistant subgroup must overturn the dominant norm
    by establishing an opinion leader that upholds a different norm, and the most likely individual to do this
    is the opinion leader in the resistant economy.

    \item \textbf{Optimization:} the optimization process is identical to before.
\end{itemize}

\section{Cross-Group Interaction}
In section's 4 and 5 we have defined the isolated social reward economies for the dominant and resistant groups respectively.
Now we wish to modify these definitions so that we can model their interaction.

First I propose a slight modification to the theoretical reputation calculation.
Let $R_i^d(\pi_i)$ be the reputation of agent $i$ in the dominant economy then define
\[ R_i^d(\pi_i) = \begin{cases}
    \gamma^{d}_r\displaystyle\sum_{j \in G \setminus \{ i \}}\sum_{s \in S}u_j(s, \pi_i(s)) & i \in G_r \\
    \gamma^d_r\displaystyle\sum_{j \in G \setminus \{ i \}}\sum_{s \in S}u_j(s, \pi_i(s)) & i \in G \setminus G_r \\
\end{cases} \]
where these are the same quantities up to the scaling factors $\gamma^d_r, \gamma^d_d > 0$. These two quantities
are meant to represent the fact that, given an agent $i \in G$, whether or not it is a member of the resistant subgroup
affects its reputation in the dominant group. In the example, a member of the \#MeToo movement
might have lower reputation in the dominant group than within the movement, as their actions and beliefs are not as greatly
appreciated in society as a whole or because of a stigma against people who participate in large activist movements.
This means that a member of the general public who contributes a roughly similar amount to society as some member of the \#MeToo
movement might have a higher reputation than the member of the movement.

Now we let $R_i^r(\pi_i)$ be the reputation of agent $i$ in $G_r$, we define
\[ R_i^r(\pi_i) = \begin{cases}
    \gamma^r_r\displaystyle\sum_{j \in G_r \setminus \{ i \}}\sum_{s \in S_r}u_j(s, \pi_i(s)) & i \in G_r \\
    \gamma^r_d\displaystyle\sum_{j \in G_r}\sum_{s \in S_r}u_j(s, \pi_i(s)) & i \in G \setminus G_r \\
\end{cases} \]
and these again represent the same quantities up to scaling factors $\gamma_r^r, \gamma^r_d > 0$. These two
quantities again represent that whether an agent is a member of $G_r$ or not affects their reputation within the subgroup $G_r$.
This is evident in the \#MeToo movement example as within the movement the opinions of other member's are given a much higher value
than those of individuals outside of the movement.

This modification addresses how the theoretical reputation should be altered to better describe cross-group interaction.
Now to modify the reputation estimates, it is clear that the key structure is gossip. Within the \#MeToo movement example
it is clear the mai instance when the movement interacts with the general public is if individuals
from each group appear in the same social situation and are similar enough that they would feel comfortable exchanging information with each other.
To formalize this notion of similarity we wish to define another metric to determine whether or not two agents gossip.
We define the similarity of agents $i$ and $j$ in $G$ at time $t$
as
\[ \operatorname{sim}_d(i, j, t, s) = \begin{cases}
    0 & t = 0 \\
    (1 - \alpha)\operatorname{sim}_d(i, j, t-1, s) + \alpha(u_i(s, \pi_i(t)) - u_i(s, \pi_j(t)))^2 & t > 0 \\
\end{cases} \]
for $\alpha \in (0, 1)$.
Then define the similarity of agent $i$ and $j$ as
\[ \operatorname{sim}_d(i, j, t) = \frac{1}{|S|}\sum_{s \in S}\operatorname{sim}_d(i, j, t, s) \]
then similarly for the subgroup $G_r$ we define for $s \in S_r$
\[ \operatorname{sim}_r(i, j, t, s) = \begin{cases}
    0 & t = 0 \\
    (1 - \alpha)\operatorname{sim}_r(i, j, t-1, s) + \alpha(u_i(s, \pi_i(t)) - u_i(s, \pi_j(t)))^2 & t > 0 \\
\end{cases} \]
for $\alpha \in (0, 1)$
then define the similarity of agent $i$ and $j$ as
\[ \operatorname{sim}_r(i, j, t) = \frac{1}{|S_r|}\sum_{s \in S_r}\operatorname{sim}_d(i, j, t, s) \]
now instead of allowing gossip in each economy to occur for all active agents,
we only allow two active agents in the economy $G_d$ to gossip at time $t$
if $\operatorname{sim}_d(i, j, t) < C_d$ where $C_d > 0$ is a fixed constant. Similarly, we only allow two agents $i$ and $j$
to gossip in $G_r$ if $\operatorname{sim}_r(i, j, t) < C_r$ where $C_r > 0$ is a fixed constant.

This captures a basic intuition and observable phenomenon: individuals will only socialize
and exchange information if they feel comfortable enough with each other, or in other words
they behave similarly in most social situations (what corresponds to a low similarity score).
This also helps model how resistance actually happens in real life. Specifically, consider the \#MeToo movement
again, how does the movement influence the general public? There are large demonstrations or rallies that bring
awareness to the issues the movement stands for, but a large part of the public acceptance of the movement comes from
casual human interaction. Namely, members of the movement will interact with the general public outside of the context of the movement
and they will do so more than the people in power. This happens because most of the general public has more in common with the general public
in most social situations than those in power, hence they will tend to interact more with individuals in the movement. These small interactions
are what allow the movement to become accepted and its ideas to spread throughout the general public. Formally,
this phenomenon is modulated by the similarity score above.

\end{document}